\chapter{THE FOURIER TRANSFORM}

\begin{definition}
    For a function $f \in L^{1}(\R,\C)$, the \eax{Fourier transform} $\hat{f}$ of $f$ is a complex-valued function on $\R$ defined by
    \begin{align}
        \hat{f}(\xi) = \int_{-\infty}^{\infty} f(x) e^{-2 \pi i x \xi} \, dx, \quad \xi \in \R.
    \end{align}
\end{definition}

One can show that $x \mapsto e^{-2 \pi i x \xi}$ is measurable. Via this, the mapping $h(x) = f(x) e^{-2 \pi i x \xi}$ is measurable, and dominated by $\abs{f}$, which is integrable. Hence, the Fourier transform is well-defined. Also, we can show that it is continuous on $\R$:
\begin{align}
    \hat{f}(\xi+h_{n})-\hat{f}(\xi) = \int_{-\infty}^{\infty} f(x) e^{-2 \pi i x \xi} (e^{-2 \pi i x h_{n}} - 1) \, dx
\end{align}
where the integrand coverges to $0$ pointwise as $h_{n} \to 0$. Thus, by the dominated convergence theorem, we have $\hat{f}(\xi+h_{n}) \to \hat{f}(\xi)$ as $h_{n} \to 0$. 
\\

The notion of continuity may be redefined by our theory of Lebesgue integration.

\begin{theorem}
    Let $X,Y$ be subintervals of $\R$, and $f:X \times Y \to \R$ be a function satisfying the following conditions.
    \begin{enumerate}
        \item For $y \in Y$, the function $f_{y}:X \to \R$ defined by $f_{y}(x) = f(x,y)$ is measurable.
        \item There is a non0negative function $g \in L^{1}(X)$ such that $\abs{f_{y}(x)} = \abs{f(x,y)} \leq g(x)$ almost everywhere on $X$ for all $y \in Y$.
        \item For a fixed $y \in Y$, $\lim_{t \to y} f(x,t) = f(x,y)$ almost everywhere on $X$. That is, $f_{x}$ is continuous almost everywhere.
    \end{enumerate}
    Then, $\int_{X} f(x,y) \, dx$ exists for each $y \in Y$, and the function $F:Y \to \R$ defined by $F(y) = \int_{X} f(x,y) \, dx$ is continuous on $Y$.
\end{theorem}
\begin{proof}
    Note that $f_{y} \in \cM(X)$ for each $y \in Y$, and $\abs{f_{y}} \leq g$ almost everywhere on $X$. Hence, by the dominated convergence theorem, $f_{y} \in L^{1}(X)$. Let $y_{n} \to y$ in $Y$, and define $G_{n}(x) \defeq f(x,y_{n})$. Then $G_{n}(x) \to f(x,y)$ almost everywhere on $X$ by the third assumption. Also, $F(y_{n}) = \int_{X} G_{n}(x) \, dx$ for each $n$. By the dominated convergence theorem, we have
    \begin{align}
        \lim_{n \to \infty} F(y_{n}) = \lim_{n \to \infty} \int_{X} G_{n}(x) \, dx = \int_{X} f(x,y) \, dx = F(y).
    \end{align}
\end{proof}

Note that the Fourier transform $\hat{f}$ is continuous on $\R$ by the above theorem, but may fail to be differentiable.

\begin{theorem}
    Let $X,Y$ be subintervals of $\R$, and $f:X \times Y \to \R$ be a function satisfying the following conditions.
    \begin{enumerate}
        \item For $y \in Y$, the function $f_{y}:X \to \R$ defined by $f_{y}(x) = f(x,y)$ is measurable, and $f_{a} \in L^{1}(X)$ for some $a \in Y$.
        \item $\frac{\partial}{\partial y}f(x,y)$ exists for each interior point $(x,y) \in X \times Y$.
        \item There exists a non-negative function $G \in L^{1}(X)$ such that $\abs{\frac{\partial f}{\partial y}(x,y)} \leq G(x)$ for all interior points $(x,y) \in X \times Y$.
    \end{enumerate}
    Then $\int_{X}f(x,y) \, dx$ exists for all $y \in Y$, and the function $F:Y \to \R$ defined by $F(y) = \int_{X} f(x,y) \, dx$ is differentiable at each interior point of $Y$, and $F'(y) = \int_{X} \frac{\partial f}{\partial y}(x,y) \, dx$.
\end{theorem}

\begin{proof}
    Since $f(x,y)-f(x,a) = (y-a)\frac{\partial}{\partial y}f(x,c)$ for some $c$ between $a$ and $y$, we have $\abs{f_{y}(x)} \leq \abs{f_{a}(x)} + \abs{y-a}G(x)$, showing that $\abs{f_{y}} \in L^{1}(X)$ for each $y \in Y$. Hence, $\int_{X} f(x,y) \, dx$ exists for all $y \in Y$. Let $y_{n} \to y$ in $Y$, and define $q_{n}(x) = \frac{f(x,y_{n})-f(x,y)}{y_{n}-y}$. Then $q_{n}(x) \to \frac{\partial f}{\partial y}(x,y)$ for each interior point $(x,y) \in X \times Y$. Also, $\abs{q_{n}(x)} \leq G(x)$ for all interior points $(x,y) \in X \times Y$. By the dominated convergence theorem, we have
    \begin{align}
        \lim_{n \to \infty} \int_{X} q_{n}(x) \, dx = \int_{X} \frac{\partial f}{\partial y}(x,y) \, dx.
    \end{align}
\end{proof}

\textit{August 27th.}
\begin{example}
    Since $\chi_{[-1,1]} \in L^{1}(\R)$, we can compute its Fourier transform:
    \begin{align}
        \hat{\chi}_{[-1,1]}(\xi) = \int_{\R} \chi_{[-1,1]}(x) e^{-2\pi i x \xi} \, dx = \int_{-1}^{1} e^{-2\pi i x \xi} \, dx = \frac{e^{-2\pi i \xi}-e^{2\pi i \xi}}{-2\pi i \xi} = \frac{\sin(2\pi \xi)}{\pi \xi} = 2 \sinc(2 \pi \xi)
    \end{align}
    where $\sinc(x) = \frac{\sin x}{x}$ for $x \neq 0$, and $\sinc(0) = 1$.    
\end{example}

For $y > 0$, define the function $F:(0,\infty) \to \R$ as
\begin{align}
    F(y) = \int_{0}^{\infty} e^{-xy} \sinc(x) \, dx.
\end{align}
Note that this integral is well-defined in the Lebesgue sense. Also, $\abs{e^{-xy} \sinc(x)} \leq e^{-xy}$ for $x,y > 0$. We claim that $F$ is continuous in its domain and vanishes at infinity. To this end, let $(y_{n})_{n \geq 1}$ be a sequence of positive numbers such that $y_{n} \to y > 0$. Then
\begin{align}
    F(y_{n}) - F(y) = \int_{0}^{\infty} e^{-xy_{n}}(1-e^{-x(y-y_{n})}) \sinc(x) \, dx
\end{align}
where $e^{-xy/2}(1+e^{-xy}) \geq e^{-xy_{n}}(1+e^{-xy_{n}})$ for sufficiently large $n$. Hence, by the dominated convergence theorem, we have $F(y_{n}) \to F(y)$ as $n \to \infty$. Also, for $y > 0$, we have
\begin{align}
    \abs{F(y)} \leq \int_{0}^{\infty} e^{-xy} \abs{\sinc(x)} \, dx \leq \int_{0}^{\infty} e^{-xy} \, dx = \frac{1}{y} \to 0 \quad \text{as } y \to \infty.
\end{align}
Now, $F'(y) = -\int_{0}^{\infty} e^{-xy} \sin(x) \, dx$ for $y > 0$. By integration by parts, and rearranging, we have
\begin{align}
    F'(y) = -\frac{1}{1+y^{2}}.
\end{align}
Thus, $F(y) = C - \tan^{-1}(y)$ for some constant $C$. Since $F(y) \to 0$ as $y \to \infty$, we have $C = \frac{\pi}{2}$. Moreover, $\lim_{y \to 0^{+}} F(y) = \int_{0}^{\infty} \sinc(x) \, dx = \frac{\pi}{2}$. If we choose subintervals $X = [0,\infty)$ and $Y = [\varepsilon, \infty)$ for some $\varepsilon > 0$, then the function $f(x,y) = e^{-xy} \sinc(x)$ satisfies the hypotheses of the above theorem. Therefore, $F$ is differentiable on $(0,\infty)$, and $F'(y) = -\int_{0}^{\infty} e^{-xy} \sin(x) \, dx$ for $y > 0$.

We also claim that $\lim_{\alpha \to \infty} \int_{0}^{\alpha} \sinc(x) \, dx = \frac{\pi}{2}$, as the improper Riemann integral. To see this, we know that
\begin{align}
    \lim_{\alpha \to \infty} \int_{0}^{\alpha} e^{-xy}\sinc(x) \, dx = \frac{\pi}{2} - \tan^{-1}(y) \quad \text{for } y > 0.
\end{align}
For $\varepsilon > 0$, choose $\beta > 0$ such that
\begin{align}
    \frac{5}{\beta} < \frac{\varepsilon}{3}.
\end{align}
By integration by parts, for all $\alpha > \beta$ and $y > 0$, we have
\begin{align}
    \abs{\int_{\beta}^{\alpha} \sinc(x) \, dx} &\leq \frac{2}{\beta}, \\
    \abs{\int_{\beta}^{\infty} e^{-xy}\sinc(x) \, dx} &\leq \frac{3}{\beta}.
\end{align}
Choose $y > 0$ sufficiently small such that
\begin{align}
    \abs{F(y)-\frac{\pi}{2}} < \frac{\varepsilon}{3}
    \quad \text{and} \quad
    \int_{0}^{\beta} (1-e^{-xy})\abs{\sinc(x)} \, dx < \frac{\varepsilon}{3}.
\end{align}
The second condition is possible since $e^{-xy} \to 1$ uniformly on $[0,\beta]$ as $y \to 0^{+}$. Hence, for all $\alpha > \beta$, we have
\begin{align}
    \abs{\int_{0}^{\alpha} \sinc(x) \, dx - \frac{\pi}{2}}
    &\leq \int_{0}^{\beta} (1-e^{-xy})\abs{\sinc(x)} \, dx
    + \abs{\int_{\beta}^{\alpha} \sinc(x) \, dx}
    + \abs{\int_{\beta}^{\infty} e^{-xy}\sinc(x) \, dx}
    + \abs{F(y)-\frac{\pi}{2}} \notag \\
    &< \varepsilon.
\end{align}
Therefore,
\begin{align}
    \lim_{\alpha \to \infty} \int_{0}^{\alpha} \sinc(x) \, dx = \frac{\pi}{2}.
\end{align}

\section{Convolution}
\textit{August 31st.}

\begin{definition}
    For $f,g \in L^{1}(\R,\C)$, the \eax{convolution} $f * g$ of $f$ and $g$ is a complex-valued function on $\R$ defined by
    \begin{align}
        (f * g)(x) = \int_{-\infty}^{\infty} f(x-t)g(t) \, dt, \quad x \in \R.
    \end{align}
\end{definition}
\noindent For example, if $g(x) = \frac{1}{2 \varepsilon} \chi_{[-\varepsilon,\varepsilon]}(x)$ for some $\varepsilon > 0$, then
\begin{align}
    (f * g)(x) = \frac{1}{2 \varepsilon} \int_{x-\varepsilon}^{x+\varepsilon} f(t) \, dt.
\end{align}
If $f$ is continuous at $x$, then $(f * g)(x) \to f(x)$ as $\varepsilon \to 0^{+}$.

\begin{proposition}
    The following hold true for $f,g,h \in L^{1}(\R,\C)$.
    \begin{enumerate}
        \item Commutativity: $f * g = g * f$.
        \item Associativity: $f * (g * h) = (f * g) * h$.
        \item If $g$ is smooth, then so is $f * g$.
        \item If $g$ is a polynomial on some interval, then so is $f * g$ on some interval.
    \end{enumerate}
\end{proposition}

Also, we have the following result for the Fourier transform of a convolution.
\begin{proposition}
    For $f,g \in L^{1}(\R,\C)$,
    \begin{align}
        \widehat{f * g} = \hat{f} \cdot \hat{g}.
    \end{align}
\end{proposition}
\begin{proof}
    Note that $f * g \in L^{1}(\R,\C)$, and
    \begin{align}
        \widehat{f * g}(\xi) = \int_{-\infty}^{\infty} (f * g)(x) e^{-2\pi i x \xi} \, dx = \int_{-\infty}^{\infty} \int_{-\infty}^{\infty} f(x-t)g(t) e^{-2\pi i x \xi} \, dt \, dx.
    \end{align}
    By Fubini's theorem, we can interchange the order of integration:
    \begin{align}
        \widehat{f * g}(\xi) = \int_{-\infty}^{\infty} g(t) \left( \int_{-\infty}^{\infty} f(x-t) e^{-2\pi i x \xi} \, dx \right) dt.
    \end{align}
    Let $u = x - t$, then $du = dx$, and we have
    \begin{align}
        \widehat{f * g}(\xi) = \int_{-\infty}^{\infty} g(t) e^{-2\pi i t \xi} \left( \int_{-\infty}^{\infty} f(u) e^{-2\pi i u \xi} \, du \right) dt = \hat{f}(\xi) \cdot \hat{g}(\xi).
    \end{align}
\end{proof}

\textit{September 17th.}

\begin{theorem}
    Let $X,Y$ be subintervals of $\R$, and suppose $k:X \times Y \to \R$ is continuous and bounded. With $f \in L^{1}(X)$ and $g \in L^{1}(Y)$, the following hold true.
    \begin{enumerate}
        \item For $y \in Y$, $\int_{X}f(x)k(x,y)dx$ exists and $F(y) = \int_{X}f(x)k(x,y)dx$ is continuous on $Y$.
        \item For $x \in X$, $\int_{Y}g(y)k(x,y)dy$ exists and $G(x) = \int_{Y}g(y)k(x,y)dy$ is continuous on $X$.
        \item The integrals $\int_{X}f(x)G(x)dx$ and $\int_{Y}g(y)F(y)dy$ exist, and
        \begin{align}
            \int_{Y}g(y)\left(\int_{X}f(x)k(x,y) \, dx\right)dy = \int_{Y} g(y)F(y) \, dy = \int_{X} f(x)G(x) \, dx = \int_{X}f(x)\left( \int_{Y} g(y)k(x,y) \, dy \right) dx.
        \end{align}
    \end{enumerate}
\end{theorem}

\begin{proof}
    Let $M = \sup_{X \times Y} \abs{k} < \infty$.

    For the first statement, fix $y \in Y$ and let $f_{y}(x) = f(x)k(x,y)$. Since $f$ is measurable and $x \mapsto k(x,y)$ is continuous, $f_{y}$ is measurable, and $\abs{f_{y}(x)} \leq M\abs{f(x)}$ with $M\abs{f} \in L^{1}(X)$. By the dominated convergence theorem, $f_{y} \in L^{1}(X)$, so $F(y) = \int_{X} f(x)k(x,y) \, dx$ exists. Now let $y_{n} \to y$ in $Y$. Then
    \begin{align}
        \abs{F(y_{n})-F(y)} \leq \int_{X} \abs{f(x)} \abs{k(x,y_{n})-k(x,y)} \, dx
    \end{align}
    where the integrand converges to $0$ pointwise as $n \to \infty$ by the continuity of $k$, and is dominated by $2M\abs{f(x)} \in L^{1}(X)$. Hence, by the dominated convergence theorem, we have $F(y_{n}) \to F(y)$ as $n \to \infty$, so $F$ is continuous on $Y$. The second statement follows by interchanging the roles of $(X,f)$ and $(Y,g)$.

    For the third statement, note that $G$ is continuous by the second statement, with $\abs{G(x)} \leq M\norm{g}_{1}$. Thus $fG$ is measurable and $\abs{f(x)G(x)} \leq M\norm{g}_{1}\abs{f(x)}$, so $fG \in L^{1}(X)$; likewise $\abs{g(y)F(y)} \leq M\norm{f}_{1}\abs{g(y)}$, so $gF \in L^{1}(Y)$. Both integrals therefore exist, and it remains to prove that $\int_{X} f(x)G(x) \, dx = \int_{Y} g(y)F(y) \, dy$. For $\phi \in L^{1}(X)$ and $\psi \in L^{1}(Y)$, write
    \begin{align}
        I(\phi,\psi) = \int_{X} \phi(x)\left(\int_{Y} \psi(y)k(x,y) \, dy\right) dx, \qquad J(\phi,\psi) = \int_{Y} \psi(y)\left(\int_{X} \phi(x)k(x,y) \, dx\right) dy,
    \end{align}
    so that the desired identity is $I(f,g) = J(f,g)$. Both $I$ and $J$ are bilinear, and since the inner integrals are bounded by $M\norm{\psi}_{1}$ and $M\norm{\phi}_{1}$ respectively,
    \begin{align}
        \abs{I(\phi,\psi)} \leq M\norm{\phi}_{1}\norm{\psi}_{1}, \qquad \abs{J(\phi,\psi)} \leq M\norm{\phi}_{1}\norm{\psi}_{1}.
    \end{align}
    If $\phi = \chi_{A}$ and $\psi = \chi_{B}$ for compact intervals $A \subseteq X$ and $B \subseteq Y$, then $k$ is continuous on the compact rectangle $A \times B$, and the theory of the multivariable Riemann integral gives
    \begin{align}
        I(\chi_{A},\chi_{B}) = \int_{A}\left(\int_{B} k(x,y) \, dy\right) dx = \int_{B}\left(\int_{A} k(x,y) \, dx\right) dy = J(\chi_{A},\chi_{B}).
    \end{align}
    By bilinearity, $I(s,t) = J(s,t)$ for all step functions $s$ on $X$ and $t$ on $Y$.

    Now let $\varepsilon > 0$, and choose step functions $s$ and $t$ with $\int_{X} \abs{f-s} \, dx < \varepsilon$ and $\int_{Y} \abs{g-t} \, dy < \varepsilon$, so that $\norm{s}_{1} \leq \norm{f}_{1} + \varepsilon$. Using bilinearity together with the bounds above,
    \begin{align}
        \abs{I(f,g)-I(s,t)} = \abs{I(f-s,g) + I(s,g-t)} \leq M\norm{g}_{1}\varepsilon + M(\norm{f}_{1}+\varepsilon)\varepsilon,
    \end{align}
    and identically $\abs{J(f,g)-J(s,t)} \leq M\norm{g}_{1}\varepsilon + M(\norm{f}_{1}+\varepsilon)\varepsilon$. Since $I(s,t) = J(s,t)$, it follows that
    \begin{align}
        \abs{I(f,g)-J(f,g)} \leq \abs{I(f,g)-I(s,t)} + \abs{J(s,t)-J(f,g)} \leq 2M\varepsilon\left(\norm{f}_{1}+\norm{g}_{1}+\varepsilon\right).
    \end{align}
    Letting $\varepsilon \to 0^{+}$ yields $I(f,g) = J(f,g)$, as required.
\end{proof}

We show more properties of the convolution.

\begin{proposition}
    Let $f,g \in L^{1}(\R)$.
    \begin{enumerate}
        \item If $f,g \in C_{c}(\R)$, then $f \ast g \in C_{c}(\R)$.
        \item If $f \in L^{1}(\R)$ and $g \in C_{c}(\R)$, then $f \ast g \in C(\R)$.
        \item Let $\tau_{p}f(x) = f(x-p)$ for $p \in \R$. Then $\tau_{p}(f \ast g) = (\tau_{p}f) \ast g = f \ast (\tau_{p}g)$.
        \item $(f \ast g)' = f' \ast g = f \ast g'$.
        \item Let $\varphi \in C_{c}^{\infty}(\R)$ be such that $\int_{\R} \varphi(x) dx = 1$, and $\varphi$ is an even non-negative function that is monotonically decreasing on $[0,\infty)$. Then, with $\varphi_{\varepsilon}(x) = \frac{1}{\varepsilon} \varphi(\frac{x}{\varepsilon})$, we have $\lim_{\varepsilon \to 0^{+}} f \ast \varphi_{\varepsilon} = f$ in $L^{1}(\R)$.
        \item If, additionally, $f \in C_{c}(\R)$, then $\lim_{\varepsilon \to 0^{+}} f \ast \varphi_{\varepsilon} = f$ uniformly on $\R$.
    \end{enumerate}
\end{proposition}
\begin{proof}
    \begin{enumerate}
        \item Since the integrand of $(f \ast g)(x) = \int_{\R} f(x-y)g(y) \, dy$ is nonzero only when $y \in \supp(g)$ and $x-y \in \supp(f)$, any $x$ with $(f \ast g)(x) \neq 0$ satisfies $x = (x-y)+y \in \supp(f) + \supp(g)$. As $\supp(f)$ and $\supp(g)$ are compact, so is their sum $\supp(f) + \supp(g)$, being the image of the compact set $\supp(f) \times \supp(g)$ under addition; hence $\supp(f \ast g) \subseteq \supp(f) + \supp(g)$ is compact. Moreover $f \in C_{c}(\R) \subseteq L^{1}(\R)$ and $g \in C_{c}(\R)$, so $f \ast g$ is continuous by (2). Therefore $f \ast g \in C_{c}(\R)$.

        \item Since $g \in C_{c}(\R)$ is bounded and uniformly continuous, commutativity gives $(f \ast g)(x) = \int_{\R} f(y)g(x-y) \, dy$, so for $h \in \R$,
        \begin{align}
            \abs{(f \ast g)(x+h)-(f \ast g)(x)} \leq \int_{\R} \abs{f(y)} \abs{g(x+h-y)-g(x-y)} \, dy \leq \norm{f}_{1} \sup_{z \in \R} \abs{g(z+h)-g(z)}.
        \end{align}
        By the uniform continuity of $g$, the supremum tends to $0$ as $h \to 0$, uniformly in $x$. Hence $f \ast g$ is (uniformly) continuous on $\R$.

        \item Directly, and using the substitution $u = y-p$ where needed,
        \begin{align}
            (\tau_{p}f \ast g)(x) &= \int_{\R} f(x-y-p)g(y) \, dy = (f \ast g)(x-p) = \tau_{p}(f \ast g)(x), \\
            (f \ast \tau_{p}g)(x) &= \int_{\R} f(x-y)g(y-p) \, dy = \int_{\R} f(x-p-u)g(u) \, du = (f \ast g)(x-p).
        \end{align}
        Hence $\tau_{p}(f \ast g) = (\tau_{p}f) \ast g = f \ast (\tau_{p}g)$.

        \item Suppose $g \in C_{c}^{1}(\R)$. Writing $(f \ast g)(x) = \int_{\R} f(y)g(x-y) \, dy$, the integrand has $x$-derivative $f(y)g'(x-y)$, which is dominated by $\norm{g'}_{\infty}\abs{f(y)} \in L^{1}(\R)$. By differentiation under the integral sign,
        \begin{align}
            (f \ast g)'(x) = \int_{\R} f(y)g'(x-y) \, dy = (f \ast g')(x).
        \end{align}
        Symmetrically, if $f \in C_{c}^{1}(\R)$, then writing $(f \ast g)(x) = \int_{\R} f(x-y)g(y) \, dy$ and differentiating under the integral sign gives $(f \ast g)'(x) = \int_{\R} f'(x-y)g(y) \, dy = (f' \ast g)(x)$. In particular, when both derivatives exist, $(f \ast g)' = f' \ast g = f \ast g'$.

        \item Since $\int_{\R} \varphi_{\varepsilon}(y) \, dy = \int_{\R} \varphi(y) \, dy = 1$ for every $\varepsilon > 0$, we may write
        \begin{align}
            (f \ast \varphi_{\varepsilon})(x) - f(x) = \int_{\R} \left(f(x-y)-f(x)\right)\varphi_{\varepsilon}(y) \, dy.
        \end{align}
        Taking the $L^{1}$-norm and using Tonelli's theorem to interchange the integrals,
        \begin{align}
            \norm{f \ast \varphi_{\varepsilon} - f}_{1} \leq \int_{\R} \varphi_{\varepsilon}(y) \left(\int_{\R} \abs{f(x-y)-f(x)} \, dx\right) dy = \int_{\R} \varphi_{\varepsilon}(y)\, \omega(y) \, dy,
        \end{align}
        where $\omega(y) = \norm{\tau_{y}f - f}_{1}$. Recall that translation is continuous in $L^{1}(\R)$, that is, $\omega(y) \to 0$ as $y \to 0$: given $\eta > 0$, choose $\psi \in C_{c}(\R)$ with $\norm{f-\psi}_{1} < \eta$, so that
        \begin{align}
            \omega(y) \leq \norm{\tau_{y}f - \tau_{y}\psi}_{1} + \norm{\tau_{y}\psi - \psi}_{1} + \norm{\psi - f}_{1} = 2\norm{f-\psi}_{1} + \norm{\tau_{y}\psi - \psi}_{1},
        \end{align}
        and $\norm{\tau_{y}\psi - \psi}_{1} \to 0$ as $y \to 0$ by the uniform continuity of $\psi$; also $\omega(y) \leq 2\norm{f}_{1}$ for all $y$. Now fix $\eta > 0$ and choose $\delta > 0$ with $\omega(y) < \eta$ for $\abs{y} < \delta$. Then
        \begin{align}
            \norm{f \ast \varphi_{\varepsilon} - f}_{1} \leq \int_{\abs{y} < \delta} \varphi_{\varepsilon}(y)\, \omega(y) \, dy + \int_{\abs{y} \geq \delta} \varphi_{\varepsilon}(y)\, \omega(y) \, dy \leq \eta + 2\norm{f}_{1} \int_{\abs{y} \geq \delta} \varphi_{\varepsilon}(y) \, dy.
        \end{align}
        Finally, $\int_{\abs{y} \geq \delta} \varphi_{\varepsilon}(y) \, dy = \int_{\abs{u} \geq \delta/\varepsilon} \varphi(u) \, du \to 0$ as $\varepsilon \to 0^{+}$, since $\varphi \in L^{1}(\R)$. Hence we have the bound $\limsup_{\varepsilon \to 0^{+}} \norm{f \ast \varphi_{\varepsilon} - f}_{1} \leq \eta$, and as $\eta > 0$ was arbitrary, $\lim_{\varepsilon \to 0^{+}} f \ast \varphi_{\varepsilon} = f$ in $L^{1}(\R)$.
    \end{enumerate}
\end{proof}

\textit{September 21st.}

\begin{theorem}
    If $f,g \in L^{1}(\R)$, then $f \ast g \in L^{1}(\R)$, and $\norm{f \ast g}_{1} \leq \norm{f}_{1}\norm{g}_{1}$. Further, if $f,g \geq 0$, then $\norm{f \ast g}_{1} = \norm{f}_{1}\norm{g}_{1}$.
\end{theorem}
\begin{proof}
    First, we show for indicator functions $\chi_{I}$ and $\chi_{J}$. We have
    \begin{align}
        (\chi_{I} \ast \chi_{J})(x) = \int_{\R} \chi_{I}(x-y)\chi_{J}(y) \, dy = \int_{\R} \chi_{(x-I)\cap J}(y) \, dy = \abs{(x-I) \cap J},
    \end{align}
    so that
    \begin{align}
        \norm{\chi_{I} \ast \chi_{J}}_{1} = \int_{\R} \abs{(x-I) \cap J} \, dx = \int_{\R} \int_{\R} \chi_{(x-I) \cap J}(y) \, dy \, dx = \int_{\R} \int_{\R} \chi_{I}(x-y)\chi_{J}(y) \, dx \, dy = \abs{I}\abs{J}.
    \end{align}
    Now let $s = \sum_{i=1}^{p} c_{i}\chi_{I_{i}}$ and $t = \sum_{j=1}^{q} d_{j}\chi_{J_{j}}$ be step functions. Then
    \begin{align}
        (s \ast t)(x) = \sum_{i,j} c_{i}d_{j}\abs{(x-I_{i}) \cap J_{j}} \implies \norm{s \ast t}_{1} = \sum_{i,j} c_{i}d_{j}\abs{I_{i}}\abs{J_{j}} \leq \norm{s}_{1}\norm{t}_{1}.
    \end{align}
    For general $f,g \in L^{1}(\R)$, we can approximate them by step functions; let $f,g \geq 0$ with $f,g \in L^{1}(\R)$. Choose step functions $s_{n} \uparrow f$ and $t_{n} \uparrow g$ pointwise almost everywhere. Then define
    \begin{align}
        \psi_{n}(x) \defeq (s_{n} \ast t_{n})(x) = \int_{\R} s_{n}(x-y)t_{n}(y) \, dy.
    \end{align}
    Since $s_{n} \uparrow f$, it converges outside a null set $N_{f}$. Thus, $(y \mapsto s_{n}(x-y)) \uparrow (y \mapsto f(x-y))$ outside $x-N_{f}$. Similarly, $(y \mapsto t_{n}(y)) \uparrow (y \mapsto g(y))$ outside $N_{g}$. So,
    \begin{align}
        \abs{\int_{\R}\psi_{n}(x) \, dx} = \abs{\int_{\R} s_{n}(x) \, dx}\abs{\int_{\R} t_{n}(x) \, dx} \leq \norm{s_{n}}_{1}\norm{t_{n}}_{1} \leq \norm{f}_{1}\norm{g}_{1}.
    \end{align}
    Thus, by the monotone convergence theorem, we have that $\psi \defeq \lim_{n \to \infty} \psi_{n}$ exists with
    \begin{align}
        \int_{\R} \psi(x) \, dx = \lim_{n \to \infty} \int_{\R} \psi_{n}(x) \, dx = \lim_{n \to \infty} \norm{s_{n}}_{1}\norm{t_{n}}_{1} = \norm{f}_{1}\norm{g}_{1}.
    \end{align}
    It remains to show that $\psi = f \ast g$ almost everywhere. For $x \in \R$, we have
    \begin{align}
        \psi(x) = \lim_{n \to \infty} \psi_{n}(x) = \lim_{n \to \infty} \int_{\R} s_{n}(x-y)t_{n}(y) \, dy = \int_{\R} f(x-y)g(y) \, dy = (f \ast g)(x).
    \end{align}
    Therefore, $\psi = f \ast g$ almost everywhere. Finally, for general $f,g \in L^{1}(\R)$, we can write $f = f_{+}-f_{-}$ and $g = g_{+}-g_{-}$, where $f_{\pm},g_{\pm} \geq 0$ are in $L^{1}(\R)$. Then
    \begin{align}
        f \ast g = (f_{+}-f_{-}) \ast (g_{+}-g_{-}) = f_{+} \ast g_{+} - f_{+} \ast g_{-} - f_{-} \ast g_{+} + f_{-} \ast g_{-},
    \end{align}
    and the result follows.
\end{proof}

\begin{theorem}[The \eax{Weierstrass approximation theorem}]
    The space of polynomial functions on $[0,1]$ is dense in $(C[0,1],\norm{\cdot}_{\infty})$.
\end{theorem}
\begin{proof}
    Let $f \in C[0,1]$. Without the loss of generality, we may assume $f(0) = f(1) = 0$, since we can set $g(x) = f(x)-f(0)-x(f(1)-f(0))$ and approximate $g$ instead. Extend $f$ to a continuous function on $\R$ by setting $f(x) = 0$ for $x \notin [0,1]$, so that $f \in C_{c}(\R) \subseteq L^{1}(\R)$. Let $Q_{n}$ be the function defined as $Q_{n}(x) = c_{n}(1-x^{2})^{n}$ for $x \in [-1,1]$ and $Q_{n}(x) = 0$ for $x \notin [-1,1]$, where $c_{n}$ is a normalization constant. Then $Q_{n} \in C_{c}(\R)$ is non-negative, even, and monotonically decreasing on $[0,\infty)$, with $\int_{\R} Q_{n}(x) \, dx = 1$. Define $P_{n}(x) = f \ast Q_{n}(x) = \int_{0}^{1} f(t)Q_{n}(t-x) \, dt$ for $x \in [0,1]$. 
\end{proof}
(INCOMPLETE)

\begin{definition}    
    Let $E$ be a set. A collection of functions $\cA \subseteq \C^{E}$ is said to be an \eax{algebra} if the following hold.
    \begin{enumerate}
        \item $f+g \in \cA$ for all $f,g \in \cA$.
        \item $fg \in \cA$ for all $f,g \in \cA$.
        \item $\lambda f \in \cA$ for all $f \in \cA$ and $\lambda \in \C$.
    \end{enumerate}
    $\cA$ is said to be \eax{uniformly closed} if $(\cA,\norm{\cdot}_{\infty})$ is complete.
\end{definition}
Examples include the set of polynomials on $[0,1]$, the set of trigonometric polynomials on $[0,2\pi]$, and the set of continuous functions on a compact interval $[a,b]$.
\begin{definition}
    Let $\cA \subseteq \C^{E}$ be an algebra.
    \begin{enumerate}
        \item $\cA$ is said to separate points on $E$ if for any $x,y \in E$ with $x \neq y$, there exists $f \in \cA$ such that $f(x) \neq f(y)$.
        \item For each $x \in E$, there exists $g \in \cA$ such that $g(x) \neq 0$; in this case, we say that $\cA$ vanishes nowhere on $E$.
    \end{enumerate}
\end{definition}

\textit{September 24th.}

\begin{theorem}
    Let $\cA \subseteq \C^{E}$ be an algebra over $E$, with $\abs{E} \geq 2$, such that $\cA$ separate points on $E$ and vanishes nowhere on $E$. Take $x_{1} \neq x_{2}$ on $E$, and $c_{1},c_{2} \in \C$. Then there exists $f \in \cA$ such that $f(x_{1}) = c_{1}$ and $f(x_{2}) = c_{2}$.
\end{theorem}
\begin{proof}
    Note that there exist $g,h,k \in \cA$ such that $g(x_{1}) \neq g(x_{2})$, $h(x_{1}) \neq 0$, and $k(x_{2}) \neq 0$. Then set $u = gk-g(x_{1})k$ and $v = gh-g(x_{2})h$, so that $u,v \in \cA$. We have $u(x_{1}) = v(x_{2}) = 0$, and $u(x_{2}) \neq 0$ and $v(x_{1}) \neq 0$. Now define
    \begin{align}
        f = c_{1}\frac{v}{v(x_{1})} + c_{2}\frac{u}{u(x_{2})} \in \cA.
    \end{align}
    Then $f(x_{1}) = c_{1}$ and $f(x_{2}) = c_{2}$.
\end{proof}

\begin{theorem}[The \eax{discrete Stone-Weierstrass theorem}]
    Let $\cA \subseteq \R^{E}$ be an algebra of real-valued functions on a finite set $E = \{x_{1},\ldots,x_{n}\}$. If $\cA$ separates points on $E$ and vanishes nowhere on $E$, then $\cA = \R^{E}$.
\end{theorem}
\begin{proof}
    Let $f \in \R^{E}$. For each $i = 1,\ldots,n$, there exists $g_{i} \in \cA$ such that $g_{i}(x_{i}) = f(x_{i})$ and $g_{i}(x_{j}) = 0$ for all $j \neq i$. Then $f = g_{1} + \cdots + g_{n} \in \cA$.
\end{proof}

\begin{theorem}[The \eax{Stone-Weierstrass theorem}]
    Let $\cA \subseteq C(K)$ be an algebra of real-valued continuous functions on a compact Hausdorff space $K$. If $\cA$ separates points on $K$ and vanishes nowhere on $K$, then $\cA$ is dense in $(C(K),\norm{\cdot}_{\infty})$.
\end{theorem}
\begin{proof}
    We show that $\overline{\cA} = C(K)$, where $\overline{\cA}$ is the uniform closure of $\cA$. Note that $\overline{\cA}$ is an algebra, separates points on $K$, and vanishes nowhere on $K$. If $f \in \overline{\cA}$, then $\abs{f} \in \overline{\cA}$; consider the map $x \mapsto \abs{x}$ on $[-a,a]$ for $a > 0$. By the Weierstrass approximation theorem, there is a sequence of polynomial functions $p_{n}$ on $[-a,a]$ such that $p_{n}(x) \to \abs{x}$ on $[-a,a]$ uniformly. For every $\eps > 0$, there exists a polynomial function $p$ on $[-a,a]$ such that $\abs{p(x)-\abs{x}} < \eps$ for all $x \in [-a,a]$. Thus, set $g = p(f) \in \overline{\cA}$, and $\norm{g-\abs{f}}_{\infty} < \eps$. Hence $\abs{f} \in \overline{\cA}$.

    By the above, if $f,g$ are in $\overline{\cA}$, then so are $\max\{f,g\}$ and $\min\{f,g\}$. Indeed, $\max\{f,g\} = \frac{1}{2}(f+g+\abs{f-g})$ and $\min\{f,g\} = \frac{1}{2}(f+g-\abs{f-g})$. We claim that if $f \in C(K)$, $x \in K$, and $\eps > 0$, then there exists $g_{x} \in \overline{\cA}$ such that $g_{x}(x) = f(x)$, and $g_{x}(t) > f(t)-\eps$ for all $t \in K$; to this end, let $y \in K$ with $y \neq x$. By the hypothesis, there exists $h_{y} \in \overline{\cA}$ such that $h_{y}(x) = f(x)$ and $h_{y}(y) = f(y)$. By continuity, there is a neighbourhood $J_{y}$ of $x$ such that $h_{y}(t) > f(t)-\eps$ for all $t \in J_{y}$. Since $y$ was arbitrary, the collection $\{J_{y}\}_{y \in K \setminus \{x\}}$ is an open cover of $K$. $K$ being compact, there exists a finite subcover $\{J_{y_{1}},\ldots,J_{y_{n}}\}$, and we can set $g_{x} = \max\{h_{y_{1}},\ldots,h_{y_{n}}\} \in \overline{\cA}$. On $K$, we have $g_{x}(t) > f(t)-\eps$ for all $t \in K$, and $g_{x}(x) = f(x)$.

    Via the above, we can show that given $f \in C(K)$ and $\eps > 0$, there exists $h \in \overline{\cA}$ such that $\norm{f-h}_{\infty} < \eps.$ Indeed, for $x \in K$, there exists $g_{x} \in \overline{\cA}$ such that $g_{x}(x) = f(x)$ and $g_{x}(t) > f(t)-\eps$ for all $t \in K$. By continuity, again, there is a neighbourhood $N_{x}$ of $x$ such that $g_{x}(t) < f(t) + \eps$ for all $t \in N_{x}$. The collection $\{N_{x}\}_{x \in K}$ is an open cover of $K$, and by compactness, there exists a finite subcover $\{N_{x_{1}},\ldots,N_{x_{m}}\}$. Set $h = \min\{g_{x_{1}},\ldots,g_{x_{m}}\} \in \overline{\cA}$. Then for all $t \in K$, we have $f(t)-\eps < h(t) < f(t)+\eps$, so that $\norm{f-h}_{\infty} < \eps$. Therefore $\overline{\cA} = C(K)$.
\end{proof}

An algebra $\cA \subseteq C(K;\C)$ is said to be \eax{self-adjoint} if $f \in \cA$ implies $\overline{f} \in \cA$, where $\overline{f}(x) = \overline{f(x)}$ for all $x \in K$. The following is a generalization of the Stone-Weierstrass theorem to complex-valued functions.

\begin{theorem}
    Suppose $\cA \subseteq C(K;\C)$ is a self-adjoint algebra of complex-valued continuous functions on a compact Hausdorff space $K$, such that $\cA$ separates points on $K$ and vanishes nowhere on $K$. Then $\cA$ is dense in $(C(K;\C),\norm{\cdot}_{\infty})$.
\end{theorem}
\begin{proof}
    Let $\cA_{\R} = \{\Re(f), \Im(f) \mid f \in \cA\}$, the real and imaginary parts of functions in $\cA$. Let $h,k \in \cA_{\R}$ be such that $h = \Re(f)$ and $k = \Im(g) = \Re(-ig)$ for some $f,g \in \cA$. Then $h+k = \Re(f-ig) \in \cA_{\R}$, and $\lambda h = \Re(\lambda f) \in \cA_{\R}$ for $\lambda \in \R$. Also, $hk = \frac{1}{4}(\abs{f+g}^{2}-\abs{f-g}^{2}) \in \cA_{\R}$, since $\abs{f}^{2} = f\overline{f} \in \cA$ by self-adjointness. Hence $\cA_{\R}$ is an algebra of real-valued continuous functions on $K$. It separates points on $K$ and vanishes nowhere on $K$, since if $x,y \in K$ with $x \neq y$, there exists $f \in \cA$ such that $f(x) \neq f(y)$, so that either $\Re(f)(x) \neq \Re(f)(y)$ or $\Im(f)(x) \neq \Im(f)(y)$; also, for each $x \in K$, there exists $f \in \cA$ such that $f(x) \neq 0$, so that either $\Re(f)(x) \neq 0$ or $\Im(f)(x) \neq 0$. By the Stone-Weierstrass theorem, $\cA_{\R}$ is dense in $(C(K;\R),\norm{\cdot}_{\infty})$. Since every function in $\cA$ can be written as a linear combination of two functions in $\cA_{\R}$, it follows that $\cA$ is dense in $(C(K;\C),\norm{\cdot}_{\infty})$.
\end{proof}

\begin{example}
    Let $K = [0,1]$, and let $\cA$ be the set
    \begin{align}
        \left\{ x \mapsto \sum_{n=-N}^{N} a_{n}e^{2\pi i n x} \;\; \Bigg| \;\; N \in \N_{0}, \; a_{n} \in \C \right\}.
    \end{align}
    $\cA$, here, is termed the set of \eax{trigonometric polynomials}. Note that $f(0) = f(1)$ for all $f \in \cA$, so $\cA$ does not separate points on $K$. However, if we identify $0$ and $1$, then $\cA$ separates points on the resulting quotient space, which is homeomorphic to the unit circle. Then $\cA$ is a self-adjoint algebra of complex-valued continuous functions on the unit circle that separates points and vanishes nowhere, so by the Stone-Weierstrass theorem, $\cA$ is dense in $(C(S^{1};\C),\norm{\cdot}_{\infty})$.
\end{example}